\documentclass[12pt]{article}
\usepackage[utf8]{inputenc}
\usepackage[T1]{fontenc}
\usepackage{hyperref}
\title{Zhuang, Central Hongshuihe Language Resources for K-12 Curriculum Development}
\author{Compiled by Claude Sonnet 5.5 (claude-sonnet-5-5) for Andr\'as Kornai}
\date{}
\begin{document}
\maketitle

\section{Introduction}
Central Hongshuihe Zhuang (ISO 639-3: zch; Glottocode cent2009) is a Kra-Dai (Tai-Kadai) language of the Zhuang macrolanguage, spoken in China (Guangxi). Joshua Project gives about 1,176,000 speakers, while Global Recordings Network lists 100,000, so the sources disagree. Castro and Hansen (2010) treat it as one of three mutually unintelligible varieties, alongside Eastern Hongshuihe (zeh) and Liuqian (zlq). The search returned only general reference and catalogue pages. It found no school materials specific to this variety.

\section{Primary Educational Resources}
No primers, textbooks or school programmes specific to zch were found in search results. Standard Zhuang is taught and written in China, but nothing specific to this variety turned up.

\section{Linguistic and Reference Materials}
\begin{itemize}
\item Wikipedia overview of the Zhuang languages: \url{https://en.wikipedia.org/wiki/Zhuang_languages}
\item Wikitongues page for zch: \url{https://wikitongues.org/languages/zch/}
\item Detect Language page for zch: \url{https://www.detectlanguage.com/languages/zch}
\end{itemize}

\section{Religious and Community Resources}
\begin{itemize}
\item Joshua Project people-group profile: \url{https://joshuaproject.net/people_groups/18739}
\end{itemize}

\section{Digital and Audio Resources}
\begin{itemize}
\item Global Recordings Network page for zch. It states that no recordings are currently available in this language, and it points to a Jesus Film Project film in the related Northern (Hongshuihe) variety: \url{https://globalrecordings.net/en/language/zch}
\end{itemize}

\section{Curriculum Development Context}
Few resources specific to Central Hongshuihe Zhuang were found. Any curriculum would probably have to draw on materials for standard Zhuang and adapt them. Doing so needs local expertise, because the varieties are not all mutually intelligible. A fuller search in Chinese-language sources is recommended.
\end{document}